Marine radar is useful for obstacle detection, situational awareness, and local mapping, but robust learning requires large labeled datasets. Because collecting and labeling maritime radar logs is expensive, many pipelines rely on synthetic data. This creates a persistent sim-to-real gap: models that perform well in simulation can degrade substantially on real sensor data.

This proposal asks whether simulator parameters can be tuned automatically so synthetic training distributions better match real maritime operating regimes. We study an RL controller that updates simulator parameters to maximize held-out real-validation performance. The controller is conditioned on broad environment type (`lake`, `river`, `coast`) with the hypothesis that context-aware tuning outperforms fixed settings and uniform randomization, especially under route and clutter shift.

The method builds on domain-randomization and adaptive sim-to-real literature \cite{tobin2017domain,peng2018sim,chebotar2018simopt,ramos2019bayessim,muratore2020bayrn,openai2019rubik,josifovski2024cdr,ho2020rap,zhao2020sim2real}. Radar transfer motivation is further supported by recent radar-domain studies \cite{kim2024soft,trinh2026digitaltwin} and by public real-radar datasets that expose weather and domain shift challenges \cite{sheeny2020radiate,ouaknine2020carrada,schumann2021radarscenes,choi2024polaris,zust2023lars}. Prior work often uses fixed/random perturbations, one-shot adaptation, or calibration pipelines without explicit maritime context. Our contribution is an environment-aware RL outer loop that optimizes downstream real-task reward directly (Fig.~\ref{fig:method-overview}) and is compared against explicit capability gaps in related work (Fig.~\ref{fig:rw-matrix}).
